% -- Curso A2: Día 1 --
\documentclass[border = 0mm]{standalone} 

% -- entrada
\usepackage[utf8]{inputenc}

% -- matemáticas y física
\usepackage{amsmath,physics}

\usepackage{fontspec}
\usepackage[default]{fontsetup}

% -- color
\usepackage[dvipsnames]{xcolor}

% -- TikZ
\usepackage{tikz}
\usetikzlibrary{
	math, % para definir variables y funciones
	angles, quotes, % para dibujar ángulos
	arrows.meta, % para dibujar flechas
}

% otras definiciones
\newcommand{\pointer}[1]{\node[red] at #1 {$\times$};}

% -- Documento principal --
\begin{document}
\begin{tikzpicture}[scale = 0.8]

	\tikzmath{
		\Width = 5.6cm; \Height = 3/4*\Width;
	};

	\node[
		inner sep = 0mm,
		minimum width = \Width,
		minimum height = \Height,	
	] (N) at (0,0) {};
	
	\clip (N.north west) rectangle (N.south east);

	\begin{scope}[
		shift = {(N)},
		shift = {(2.65cm,-1.8cm)},
	]

	\tikzmath{
		% parámetros del plano
		\th = 24; \l = 5.8;
		% parámetros del tiro
		\phIni = 60; \vIni = 1.97;
		\xIni = -0.86*\l*cos(\th); \yIni = -tan(\th)*\xIni;
		\g = 1;
		% trayectoria
		function A(\ph) {
			return -(\g)/(2*(\vIni)^2*cos(\ph)^2);
		};
		function B(\ph) {
			return tan(\ph) + (\g*\xIni)/((\vIni)^2*cos(\ph)^2);
		};
		function C(\ph) {
			return \yIni - tan(\ph)*\xIni - (\g*(\xIni)^2)/(2*(\vIni)^2*cos(\ph)^2);
		};
		%
		function f(\x,\ph) {
			return A(\ph)*(\x)^2 + B(\ph)*\x + C(\ph);
		};
		% función del alcance
		\S = 2*(\vIni)^2*sin(\th + \phIni)*cos(\phIni)/(\g*cos(\th)^2);
	};	

	% definir coordenadas
	\coordinate (A) at (0,0);
	\coordinate (B) at (180 - \th:\l);
	\coordinate (C) at ({-\l*cos(\th)},0);

	% -- ángulo de inclinación del plano
	\pic[
		draw,
		"$\theta$",
		angle radius = 9mm,
		angle eccentricity = 1.2,
	] {angle = B--A--C};
	
	\draw[semithick, gray] (A) to (B) to (C) to cycle;

	% -- suelo
	\fill[gray!15]
		(C)++(-0.5,0) coordinate (tmp1)
		rectangle
		(0.5,-0.4) coordinate (tmp2);
		
	\draw[black] (C)++(-0.5,0) to (0.5,0);
	
	% -- trayectoria		
	\begin{scope}

	% máscara de recorte
	\clip (A)--(B)--++(0,2.5)-|cycle;
	
		% curva
		\draw[dashed, smooth]
			plot[domain = \xIni:0, samples = 60]
			(\x,{f(\x,\phIni)});
		
		% algunos puntos en la trayectoria
		\foreach \t in {0.3,0.6,...,3} {
			\tikzmath{
				\x = \xIni + \vIni*\t;
				\y = f(\x,\phIni);
			}
			
			\draw[
				semithick,
				gray, fill = gray!50,
				opacity = \t/3
			]
				(\x,\y) circle (1mm);
		}
		
	\end{scope}
	
	% -- descripción del p. inicial
	\coordinate (pIni) at (\xIni,\yIni);
	
	\draw
		(pIni) -- ++(1,0) coordinate (a);
		
	% velocidad inicial: vector
	\draw[very thick, -{Latex[scale = 0.75]}]
		(pIni) -- ++(\phIni:0.45*\vIni) coordinate (b);	
	\node[xshift = -2.8mm, yshift = 0.8mm] at (b) {$v_0$};
	
	% ángulo inicial
	\pic[
		draw,
		"$\phi$",
		angle radius = 4mm,
		angle eccentricity = 1.5,
	] {angle = a--pIni--b};
		
	\node[shift = {(-0.6mm,-4.5mm)}] at (pIni) {\small $(x_0,y_0)$};
	
	% alcance
	\draw[thick]
		(pIni)
		--node[below left, black] {$S$}
		++(-\th:\S) coordinate (pFin);
	
	\fill (pIni) circle (0.5mm);
	\fill (pFin) circle (0.5mm);
	
	% -- flecha de la gravedad
	\coordinate (g) at (-0.3,2.8);
	
	\fill[gray!75] (g) circle (0.5mm);
	\draw[very thick, -{Latex[scale = 0.75]}, gray!75]
		(g) --node[pos = 0.3, right, black] {$g$} ++(0,-0.8);
	
	\end{scope}
	
\end{tikzpicture}
\end{document}